\section{Conclusion}
\label{sec:conc}

We presented AgentFail, a diagnostic framework that decomposes LLM-agent failures into five stages and distinguishes silent from loud failures. Across 5{,}984 runs on five frontier models and three task suites, we found that silent-failure rates are task-dependent (86\% on synthetic traps, 4\% on UCI, 19\% on DSBench), that agents never self-correct silent failures in the standard protocol (0\% recovery), but that feedback signals enable 15.7\% recovery overall---with silent failures remaining far harder to correct (10.4\%) than loud failures (36.7\%). Oracle-guided repair matches the no-op baseline (64\%), establishing an upper bound on reparability. Harness ablation on both synthetic and UCI tasks confirms silent failures persist across three harness variants. We separate infrastructure failures (38.1\% of runs) from genuine model failures and provide an enriched manifest with 18 fields. We define three diagnostic benchmark tasks (observability classification, failure-stage classification, originating-step localization) with train/dev/test splits and baselines (silent heuristic: 80.3\% accuracy). We release all code, data, and annotated runs as a public benchmark. The most critical remaining gap is human taxonomy validation with $\geq 3$ annotators.
